\documentclass[12pt, a4paper]{article}

\usepackage[french]{babel}
\usepackage[utf8]{inputenc}
\usepackage[T1]{fontenc}
\usepackage{tgadventor}
\renewcommand*\familydefault{\sfdefault}
\usepackage{amsmath, amssymb}
\usepackage[top=2cm, bottom=2cm, left=2cm, right=2cm]{geometry}
\usepackage{enumitem}
\usepackage{multicol}
\usepackage{xcolor}
\usepackage{mdframed}
\usepackage{verbatim}
\usepackage{tabularx}
\usepackage{tkz-euclide}

\definecolor{vlmblue}{RGB}{25, 90, 160}
\definecolor{vlmgray}{RGB}{245, 245, 245}
\definecolor{vlmgreen}{RGB}{0, 130, 80}

% ── Titre de l'exercice dans un encadré ──
\newcommand{\titrexercice}[2]{%
  \begin{mdframed}[
    backgroundcolor=vlmgray,
    linecolor=black,
    linewidth=1.5pt,
    innerleftmargin=10pt,
    innerrightmargin=10pt,
    innertopmargin=6pt,
    innerbottommargin=6pt,
    skipabove=1em,
    skipbelow=0.5em
  ]
  \textbf{Exercice #1}%
  \ifx&#2&\else\hfill\textbf{#2}\fi
  \end{mdframed}
}

% ── Corps de l'exercice sans cadre ──
\newenvironment{exercice}[2]{%
  \titrexercice{#1}{#2}
  \begin{list}{}{%
    \setlength{\leftmargin}{0pt}
    \setlength{\rightmargin}{0pt}
    \setlength{\listparindent}{0pt}
    \setlength{\itemindent}{0pt}
    \setlength{\parsep}{0pt}
  }\item[]\noindent\ignorespaces
}{%
  \end{list}
  \vspace{1em}
}

% ── Solution ──
\ifdefined\AVECSOLUTIONS
  \newenvironment{solution}[1]{%
    \vspace{0.3em}
    \textbf{\textcolor{vlmgreen}{Solution #1}}
    \vspace{0.3em}\par
  }{%
    \vspace{1em}
  }
\else
  \newenvironment{solution}[1]{\comment}{\endcomment}
\fi

\pagestyle{empty}

\begin{document}
\begin{exercice}{EX-CH02-012}{Multiplications et divisions de nombres relatifs}
{\small
\bigskip
\def\arraystretch{1.7}
\begin{tabularx}{\linewidth}{*{4}{X}}
\textbf{A.} & \textbf{B.} & \textbf{C.} & \textbf{D.} \\
$(+5)\cdot(-3)= \dotfill$    & $(-45):(+5)= \dotfill$  & $(+56):(-7)= \dotfill$      & $(+27):(+9)= \dotfill$      \\
$(+6)\cdot(+12)= \dotfill$   & $(-14):(-7)= \dotfill$  & $(-10)\cdot(-11)= \dotfill$ & $(+4)\cdot(-8)= \dotfill$   \\
$(+8)\cdot(+11)= \dotfill$   & $(+90):(-10)= \dotfill$ & $(+9)\cdot(-9)= \dotfill$   & $(-63):(+7)= \dotfill$      \\
$(+42):(-7)= \dotfill$       & $(-72):(+12)= \dotfill$ & $(-12):(-4)= \dotfill$      & $(+11)\cdot(-11)= \dotfill$ \\
$(+96):(-12)= \dotfill$      & $(+6)\cdot(-4)= \dotfill$ & $(+25):(-5)= \dotfill$    & $(-49):(-7)= \dotfill$      \\
$(+48):(-8)= \dotfill$       & $(+28):(+7)= \dotfill$  & $(-5)\cdot(-6)= \dotfill$   & $(+10):(-10)= \dotfill$     \\
\end{tabularx}
}
\end{exercice}
\begin{solution}{EX-CH02-012}

\def\arraystretch{1.7}
\begin{tabularx}{\linewidth}{*{4}{X}}
\textbf{A.} & \textbf{B.} & \textbf{C.} & \textbf{D.} \\
$(+5)\cdot(-3)=-15$    & $(-45):(+5)=-9$   & $(+56):(-7)=-8$       & $(+27):(+9)=+3$      \\
$(+6)\cdot(+12)=+72$   & $(-14):(-7)=+2$   & $(-10)\cdot(-11)=+110$& $(+4)\cdot(-8)=-32$  \\
$(+8)\cdot(+11)=+88$   & $(+90):(-10)=-9$  & $(+9)\cdot(-9)=-81$   & $(-63):(+7)=-9$      \\
$(+42):(-7)=-6$        & $(-72):(+12)=-6$  & $(-12):(-4)=+3$       & $(+11)\cdot(-11)=-121$\\
$(+96):(-12)=-8$       & $(+6)\cdot(-4)=-24$ & $(+25):(-5)=-5$     & $(-49):(-7)=+7$      \\
$(+48):(-8)=-6$        & $(+28):(+7)=+4$   & $(-5)\cdot(-6)=+30$   & $(+10):(-10)=-1$     \\
\end{tabularx}

\end{solution}
\end{document}